\chapter*{Prólogo}
\addcontentsline{toc}{chapter}{Prólogo}

La modernización de sistemas energéticos rurales exige soluciones técnicamente sólidas,
pero también razonables desde el punto de vista económico y operativo. En ese cruce
entre teoría y necesidad práctica se ubica este libro, cuyo eje central es la sintonización
de controladores PID mediante métodos metaheurísticos aplicados a una microred
hidro-solar aislada de contexto andino.

El texto parte de una constatación sencilla: aunque el controlador PID sigue siendo la
herramienta de regulación más difundida en la industria, su desempeño depende de forma
decisiva del método con que se ajustan sus ganancias. En plantas no lineales, con baja
inercia o sujetas a perturbaciones severas, las reglas clásicas suelen quedarse cortas.
Por ello, esta obra explora una ruta de mejora basada en algoritmos de optimización como
PSO, GWO, Eagle Strategy y ABC, siempre con foco en la respuesta dinámica del sistema
y no en la novedad algorítmica por sí sola.

Más que presentar un inventario de técnicas, el libro propone una lectura comparativa:
cada capítulo aporta fundamentos, criterios de evaluación y elementos metodológicos para
entender por qué una sintonización puede ser más robusta, más rápida o más adecuada
para una determinada planta. Esa orientación hace que el material resulte útil tanto para
estudiantes de ingeniería como para investigadores y profesionales que buscan una base
ordenada para desarrollar sus propias pruebas en MATLAB/Simulink u otros entornos.
